\scanpage{23}\mspage{18 ter}
% @ top, continuation | text | ± εᵢ ± εⱼ
$\pm\varepsilon_i\pm\varepsilon_j$ $(1\le i<j\le5)$, et l'on a un isom canonique $E(\Pi)\simeq E$. De façon intrinsèque :

% @ upper third, a | text | les x + σx
a) les $x+\sigma x$ $(x\in E(u))$ sont tous égaux, soit $\theta\in\widehat E(C)$ [savoir $\theta=\eta-\unsure{\zeta|\xi|\xi'}_6$]

% @ upper third through fold, b | text | L'hom x ↦ x − ½θ
b) L'hom $x\longmapsto x-\frac12\theta$ de $E(u)$ dans $\frac12\widehat E(C)$ \unsure{induit|\deleted{induit}} un iso de $\mathbb Z^{E(u)}/N\simeq E(\Pi)$ \deleted{sur} \uline{dans} $\unsure{\deleted{\widehat{\vphantom E}}E(C)|\widehat E(C)}$, et les différences \deleted{successives} \inserted{mutuelles} \inserted{de $E(\Pi)$} d'éléments non liés de $E(u)$ sont \unsure{transformées|transformés} en les diff. mutuelles de ces éléments dans $\widehat E(C)$, lesquelles forment donc un syst. de \inserted{(40)} racines de type $D_5$ -- qui sont aussi les \inserted{40} racines de $E_6(C)$ \deleted{\unsure{ou}} $=E(C)$ \uline{orthogonales} à $u$.
% archive sheet number: 22
